\section{Open-loop analysis}

The open-loop analysis is performed on the linearized model. Its purpose is to characterize the plant before designing the controllers. The analysis includes stability, controllability, observability and the practical operating range.

\subsection{Stability analysis}

The open-loop poles are obtained from the eigenvalues of the state matrix $A$:
\begin{equation}
    \lambda_i = \mathrm{eig}(A).
\end{equation}

Because the mechanical dynamics of the magnetic levitation system are unstable, at least one pole is expected to be in the right half-plane. This confirms that feedback control is required to keep the ball levitated around the selected equilibrium point.

\begin{lstlisting}[style=matlabstyle,caption={Open-loop pole computation in MATLAB.}]
poles_ol = eig(A_ext);
disp(poles_ol);
\end{lstlisting}

\subsection{Controllability}

The controllability matrix is
\begin{equation}
    \mathcal{C} = \begin{bmatrix} B & AB & A^2B \end{bmatrix}.
\end{equation}

The linearized system is controllable if
\begin{equation}
    \mathrm{rank}(\mathcal{C}) = n,
\end{equation}
where $n$ is the number of states.

\begin{lstlisting}[style=matlabstyle,caption={Controllability check.}]
Mr = ctrb(A_ext, B_ext);
rank_Mr = rank(Mr);
\end{lstlisting}

\subsection{Observability}

The observability matrix is
\begin{equation}
    \mathcal{O} =
    \begin{bmatrix}
        C \\
        CA \\
        CA^2
    \end{bmatrix}.
\end{equation}

The system is observable if
\begin{equation}
    \mathrm{rank}(\mathcal{O}) = n.
\end{equation}

\begin{lstlisting}[style=matlabstyle,caption={Observability check.}]
Mo = obsv(A_ext, C_ext);
rank_Mo = rank(Mo);
\end{lstlisting}

\subsection{Operating range and constraints}

The main physical constraints are the finite travel of the ball, the maximum continuous coil current and the voltage limit of the amplifier. These constraints are relevant both for simulation and experimental implementation, because a controller that performs well in the linear model may still fail when the reference or the control effort pushes the system outside the admissible region.

The most important constraints considered in the controller design are
\begin{equation}
    0 \leq i(t) \leq i_{max},
    \qquad
    |u(t)| \leq u_{max},
    \qquad
    x_{min} \leq x(t) \leq x_{max}.
\end{equation}

The current saturation is particularly relevant when integral action is used, because actuator saturation can produce integrator windup if no protection strategy is included.
